\section{Discussion}
\label{sec:discussion}
The practical choice begins with an objective. Coverage is suitable for distributing a review surface across distinct observed endpoints; additive H allocates selection toward account-origin pair mass under the stipulated weights; P requires selecting every observed source of a nonempty endpoint. The measured tradeoffs show that a defender cannot substitute one objective's score for another's. None of these quantities estimates how many incidents a real control prevents.

The additional experiments also distinguish useful evidence from mathematical expectation. Exact-union ranking's additive optimality follows from the objective; the empirical contribution is the associated coverage and exclusion cost. The new cohort repeats that direction while changing effect sizes. A representation effect of 42.44\% in one union and 3.33\% in the other warrants reporting population and version context whenever reach is compared.

The exclusion heuristic illustrates why added complexity needs empirical justification. Reciprocal scoring is unnecessary for its original-graph P improvement, yet helps relative to the immediate variant in the new cohort. Candidate-restricted exact instances expose large empirical gaps despite frequent optimum matches. The present results support a transparent heuristic comparison, without a universal recommended cap or algorithmic guarantee.

The endpoint audit establishes a reproducible evidence-collection and review process. Its current automated findings do not settle mapping precision, execution or historical validity. Independent annotations and adjudication must precede a human-validated endpoint claim. A deployment recommendation would additionally require observed authorization paths, intervention costs, feasible control policies and outcome evaluation. These missing measurements are separate from whether the retained graph computations are correct.
